
\sourcenote{Genre: admissible rule (tier 2) --- but of the family's
third kind. \rulename{PConstExact} is admissible by reduction to
verified machinery; its siblings here are admissible with shipped
direct lemmas; \rulename{ReLU-Split} is admissible in the purest
sense: it has \emph{no proof-term constructor and no wire tag}, and a
producer elaborates it into the derivation below --- a certificate
stream never contains a \rulename{ReLU-Split} node. The unpacking
\emph{is} the rule; its soundness is its derivability --- and, since
the mechanization shipped, also a theorem: \texttt{sound\_reluSplit}
states the rule itself, at the rank-$0$ scope
Remark~\ref{rem:relusplit-scope} proves is the rule's true extent.}

\begin{brecipe}[the \rulename{ReLU-Split} elaboration]
\label{rec:relusplit-instance}
The rule, as the source catalogue states it:
\[
\inferrule*[Right=\rulename{ReLU-Split}]{
  \mathcal{T}_{+} \rhd \Phi \wedge (h \geq 0) \vdash \Psi\sub{b}{h} \\
  \mathcal{T}_{-} \rhd \Phi \wedge (h \leq 0) \vdash \Psi\sub{b}{0}
}{
  \Phi \vdash \Psi
}
\]
where $b$ marks the occurrences of the output wire $y$ of a ReLU
application $\ReLU_{\{h\}\to\{y\}} \in H$, so that $\Psi =
\Psi\sub{b}{y}$. The producer emits, in place of any such node, the
derivation:
\begin{enumerate}
\item the node's sign-split fact --- $(h \leq 0 \wedge y = 0) \vee
  (0 \leq h \wedge y = h)$, the packaged semantic contract stated in
  the \rulename{LiRPA} section's family template --- enters as a
  sub-conjunct of the premise (the sp-axiom convention);
\item $\vee$E \emph{in its context-carrying form} splits on it (see
  Remark~\ref{rem:relusplit-ctxve});
\item \emph{active branch} (context gains $0 \leq h$ and $y = h$):
  the branch context aligns with $\mathcal{T}_{+}$'s premise $\Phi
  \wedge (h \geq 0)$ by $\wedge$E, \rulename{Weaken}, and the
  orientation plumbing of the \rulename{ReLU-Above} section;
  $\mathcal{T}_{+}$ yields $\Psi\sub{b}{h}$, and \rulename{EqElim}
  along $h = y$ (wire-to-wire, the node's one shared shape)
  transports it to $\Psi\sub{b}{y} = \Psi$;
\item \emph{inactive branch} (context gains $h \leq 0$ and $y = 0$):
  symmetrically, $\mathcal{T}_{-}$ yields $\Psi\sub{b}{0}$, and
  \rulename{EqElim} along $0 = y$ transports it to $\Psi$;
\item \rulename{Cut} and \rulename{Weaken} stitch the case analysis
  back under the original $\Phi$.
\end{enumerate}
Step 4's transport has a rank-$0$ constant side, so at tensor-shaped
instances it is stopped by \rulename{EqElim}'s broadcast-shape gate,
exactly as the family's shape-scope remark records
(Remark~\ref{rem:relunn-shapescope}, \rulename{ReLU-NN} section) ---
at such instances a producer splits per scalar node
(\rulename{ReLU-Neg}'s direct node supplies the shape-general fact
$y = 0$, but transporting it into an arbitrary $\Psi$ rides the same
gate, so the per-scalar split is the general route --- and
necessarily so: Remark~\ref{rem:relusplit-scope} proves the
tensor-shaped reading admits a countermodel).
\end{brecipe}

\begin{bremark}[why the context-carrying $\vee$E]
\label{rem:relusplit-ctxve}
The disjunction being eliminated does not stand alone as the premise:
it arrives \emph{inside} the premise conjunction, next to everything
else $\Phi$ carries. The plain sequent form of $\vee$E ($\Phi_1
\vdash \Psi$ and $\Phi_2 \vdash \Psi$ give $\Phi_1 \vee \Phi_2 \vdash
\Psi$) cannot reach it --- and this is not a stylistic gap: the
$\vee$E section records the countermodel showing the context-carrying
form is strictly stronger than the printed one (no elaboration of the
former into the latter exists). The derivation above is, historically,
the essential consumer that forced that ruling; the
context-carrying form is proven sound as a satellite lemma
(\texttt{sound\_orECtx}).
\end{bremark}

\begin{bremark}[mechanization status --- the rule-level statement]
\label{rem:relusplit-mech}
The checker column remains empty for the structural reason it always
was: no certificate ever contains the rule, so there is no payload to
validate and no soundness case in the checker's recursion. The
rule-level statement this remark once recorded as follow-up work now
exists: \texttt{sound\_reluSplit} (Qtsl/ReluSplit.lean) is
the semantic lemma --- from the validity of the two branch sequents
and the linkage, the conclusion sequent is valid --- built from
exactly the vocabulary the remark predicted (the linkage extraction
\texttt{reluLinked\_facts}, the substitution predicate
\texttt{substituteInFormula} that \rulename{EqElim}'s payload already
uses, and the semantic transport \texttt{Formula.holds\_subst\_congr}
that \rulename{EqElim}'s soundness already rides). The template
$\Psi$ with its marked occurrences $b$ is an arbitrary formula over an
arbitrary marker wire, with the three substitution instances
($\Psi\sub{b}{h}$, $\Psi\sub{b}{0}$, $\Psi = \Psi\sub{b}{y}$) as
hypotheses --- no freshness condition, since the proof transports
along value equality, never syntax. ``Derivable, by inspection of
this recipe'' is thereby upgraded to ``sound, by theorem'' at every
instance the recipe itself reaches (rank-$0$ nodes ---
Remark~\ref{rem:relusplit-scope}); the stronger recorded follow-up, a
verified elaboration function producing
Recipe~\ref{rec:relusplit-instance}'s derivation, remains open.
\end{bremark}

\begin{bremark}[the rank-$0$ scope is semantic necessity]
\label{rem:relusplit-scope}
\texttt{sound\_reluSplit} hypothesizes the output wire typed at
rank~$0$, and this is the rule's true extent rather than proof
convenience. The family's shape-scope remark
(Remark~\ref{rem:relunn-shapescope}) records that the
\emph{derivation}'s constant-sided transport is stopped by
\rulename{EqElim}'s broadcast-shape gate at tensor shapes; for the
siblings the direct lemma is nonetheless shape-general, so the gate
reads there as a technicality of the derivation route. Not here:
\texttt{reluSplit\_rank0\_necessary} (same module) mechanizes a
countermodel showing the shape-general reading of
\rulename{ReLU-Split} is \emph{false}. Template $\Psi(b) = (b = 0)
\vee (b = h)$ over a shape-$[2]$ node, $\Phi = \top$: both premises
are valid ($\mathcal{T}_{+}$ through $h = h$, $\mathcal{T}_{-}$
through $0 = 0$), yet at the mixed-sign valuation $h = [1, -1]$,
$y = [1, 0]$ neither disjunct of the conclusion holds. Mixed signs
are exactly what the rule's case analysis cannot see from
sign-definite branch hypotheses --- the semantic content of ``each
split doubles the branch count'': a tensor node is $2^{n}$ scalar
cases, not $2$. Even shape $[1]$ fails, by a different leak
(broadcast-arity mixing: $\Psi(b) = (b = h) \vee (b + w \geq 0)$ with
$w$ typed at $[2]$ --- the $[1]$-shaped $y + w$ atom is unsatisfiable
by the strict-shape arithmetic while the rank-$0$ constant broadcasts;
recorded as prose in the module header). The recipe's step-4 note
(``the per-scalar split is the general route'') is thus not a
workaround for a gate but the only sound reading; producers split per
rank-$0$ node, and \texttt{sound\_reluSplit} is the rule those splits
instantiate.
\end{bremark}

\begin{bremark}[when a producer splits]
Splitting is the \emph{complete}-verification move of the
neural-network literature --- Reluplex's case
analysis~\cite{katz2017reluplex}, systematized as branch and
bound~\cite{bunel2020branch}: exact where \rulename{LiRPA}'s envelope
is lossy, at the price of one branch doubling per split neuron, so
proof size is exponential in the split count. The practical frontier
--- which neurons to split, which to relax --- is the branch-and-bound
literature's whole subject; the calculus prices the choice honestly
(a split is two subtrees, an envelope is one node) and takes no
position on it.
\end{bremark}

\begin{anchors}
the sign-split fact (step 1) & \texttt{PrimApp.SemOk} --- the packaged contract; conservativity via \texttt{relu\_projection} (\rulename{LiRPA} section)\\
context-carrying case split (step 2) & \texttt{sound\_orECtx} (satellite lemma, Admissible.lean)\\
equality transport (steps 3--4) & \texttt{sound\_eqElim}; \texttt{substituteInFormula}\\
rule-level statement (rank-$0$ scope) & \texttt{sound\_reluSplit} (ReluSplit.lean) --- Remark~\ref{rem:relusplit-mech}\\
scope necessity (mixed-sign countermodel) & \texttt{reluSplit\_rank0\_necessary}; fixture \texttt{ReluSplitScope.HSplit} --- Remark~\ref{rem:relusplit-scope}\\
verified elaboration function & \emph{none exists} --- the remaining recorded follow-up\\
\end{anchors}
